\SourceLeaf{120}
\[
\frac{m\Pi\,\frac12(n-1)}{h^n\surd\pi}\Bigl\{1-\rho\surd\Bigl(\frac2m.\Bigl(\frac{\Pi(n-\frac12).\surd\pi}{(\Pi\,\frac12(n-1))^2}-1\Bigr)\Bigr)\Bigr\}
\]
und
\[
\frac{m\Pi\,\frac12(n-1)}{h^n\surd\pi}\Bigl\{1+\rho\surd\Bigl(\frac2m.\Bigl(\frac{\Pi(n-\frac12).\surd\pi}{(\Pi\,\frac12(n-1))^2}-1\Bigr)\Bigr)\Bigr\}
\]
Setzt man also, wie oben,
\[
\frac\rho h=r
\]
so dass \(r\) den wahrscheinlichen Beobachtungsfehler vorstellt, so ist der wahrscheinlichste Werth von
\[
\rho\BareRoot[n]{\frac{S^{\mathrm n}\surd\pi}{m\Pi\,\frac12(n-1)}}
\]
offenbar \(=r\); und die wahrscheinlichen Grenzen des Werthes jener Grösse
\[
r\Bigl\{1-\frac\rho n\surd\Bigl(\frac2m.\Bigl(\frac{\Pi(n-\frac12).\surd\pi}{(\Pi\,\frac12(n-1))^2}-1\Bigr)\Bigr)\Bigr\}
\]
und
\[
r\Bigl\{1+\frac\rho n\surd\Bigl(\frac2m.\Bigl(\frac{\Pi(n-\frac12).\surd\pi}{(\Pi\,\frac12(n-1))^2}-1\Bigr)\Bigr)\Bigr\}
\]

Es ist also auch eins gegen eins zu wetten, dass \(r\) zwischen den Grenzen
\[
\begin{aligned}
&\rho\BareRoot[n]{\frac{S^{\mathrm n}\surd\pi}{m\Pi\,\frac12(n-1)}}\\[3pt]
&\quad\Bigl\{1-\frac\rho n\surd\Bigl(\frac2m.\Bigl(\frac{\Pi(n-\frac12).\surd\pi}{(\Pi\,\frac12(n-1))^2}-1\Bigr)\Bigr)\Bigr\}
\end{aligned}
\]
und
\[
\begin{aligned}
&\rho\BareRoot[n]{\frac{S^{\mathrm n}\surd\pi}{m\Pi\,\frac12(n-1)}}\\[3pt]
&\quad\Bigl\{1+\frac\rho n\surd\Bigl(\frac2m.\Bigl(\frac{\Pi(n-\frac12).\surd\pi}{(\Pi\,\frac12(n-1))^2}-1\Bigr)\Bigr)\Bigr\}
\end{aligned}
\]
liege. Für \(n=2\) sind diese Grenzen
\[
\rho\surd\frac{2S''}{m}\bigl\{1-\tfrac\rho{\surd m}\bigr\}\quad\text{und}\quad
\rho\surd\frac{2S''}{m}\bigl\{1+\tfrac\rho{\surd m}\bigr\}
\]
ganz mit den oben (Art. 4) gefundenen übereinstimmend. Allgemein hat man für ein gerades \(n\) die Grenzen
\[
\begin{aligned}
&\rho\surd2.\BareRoot[n]{\frac{S^{\mathrm n}}{m.1.3.5.7\, .\,.\,.\,.\,(n-1)}}\\[4pt]
&\quad\Bigl\{1-\frac\rho n\surd\Bigl(\frac2m.\Bigl(\frac{(n+1).(n+3)\, .\,.\,.\,(2n-1)}{1.3.5\, .\,.\,.\,(n-1)}-1\Bigr)\Bigr)\Bigr\}
\end{aligned}
\]
und
\[
\begin{aligned}
&\rho\surd2.\BareRoot[n]{\frac{S^{\mathrm n}}{m.1.3.5.7\, .\,.\,.\,.\,(n-1)}}\\[4pt]
&\quad\Bigl\{1+\frac\rho n\surd\Bigl(\frac2m.\Bigl(\frac{(n+1).(n+3)\, .\,.\,.\,(2n-1)}{1.3.5\, .\,.\,.\,(n-1)}-1\Bigr)\Bigr)\Bigr\}
\end{aligned}
\]
und für ein ungerades \(n\) folgende
\[
\begin{aligned}
&\rho\BareRoot[n]{\frac{S^{\mathrm n}\surd\pi}{m.1.2.3\, .\,.\,.\,\frac12(n-1)}}\\[4pt]
&\quad\Bigl\{1-\frac\rho n\surd\Bigl(\frac1m.\Bigl(\frac{1.3.5.7\, .\,.\,.\,(2n-1)\pi}{(2.4.6\, .\,.\,.\,(n-1))^2}-2\Bigr)\Bigr)\Bigr\}
\end{aligned}
\]
und
\[
\begin{aligned}
&\rho\BareRoot[n]{\frac{S^{\mathrm n}\surd\pi}{m.1.2.3\, .\,.\,.\,\frac12(n-1)}}\\[4pt]
&\quad\Bigl\{1+\frac\rho n\surd\Bigl(\frac1m.\Bigl(\frac{1.3.5.7\, .\,.\,.\,(2n-1)\pi}{(2.4.6\, .\,.\,.\,(n-1))^2}-2\Bigr)\Bigr)\Bigr\}
\end{aligned}
\]
\SourceLeaf{121}
\Article{6}
Ich füge noch die numerischen Werthe für die einfachsten Fälle bei:
\begin{center}Wahrscheinliche Grenzen von \(r\)\end{center}
\[
\begin{array}{r@{\quad}l}
\mathrm{I}.&0{,}8453473\quad\dfrac{S'}m.\bigl(1\mp\dfrac{0{,}5095841}{\surd m}\bigr)\\[10pt]
\mathrm{II}.&0{,}6744897\BareRoot[2]{\dfrac{S''}m}.\bigl(1\mp\dfrac{0{,}4769363}{\surd m}\bigr)\\[10pt]
\mathrm{III}.&0{,}5771897\BareRoot[3]{\dfrac{S'''}m}.\bigl(1\mp\dfrac{0{,}4971987}{\surd m}\bigr)\\[10pt]
\mathrm{IV}.&0{,}5125017\BareRoot[4]{\dfrac{S''''}m}.\bigl(1\mp\dfrac{0{,}5507186}{\surd m}\bigr)\\[10pt]
\mathrm{V}.&0{,}4655532\BareRoot[5]{\dfrac{S^{\mathrm V}}m}.\bigl(1\mp\dfrac{0{,}6355080}{\surd m}\bigr)\\[10pt]
\mathrm{VI}.&0{,}4294972\BareRoot[6]{\dfrac{S^{\mathrm{VI}}}m}.\bigl(1\mp\dfrac{0{,}7557764}{\surd m}\bigr)
\end{array}
\]

Man sieht also auch hieraus, dass die Bestimmungsart II von allen die vortheilhafteste ist. Hundert Beobachtungsfehler, nach dieser Formel behandelt, geben nemlich ein eben so zuverlässiges Resultat, wie
\begin{center}
114 nach I,\quad109 nach III,\quad133 nach IV,\quad178 nach V,\quad251 nach VI.
\end{center}

Inzwischen hat die Formel I den Vorzug der allerbequemsten Rechnung. und man mag sich daher derselben, da sie doch nicht viel weniger genau ist als II, immerhin bedienen, wenn man nicht die Summe der Quadrate der Fehler sonst schon kennt, oder zu kennen wünscht.

\Article{7}
Noch bequemer, obwohl beträchtlich weniger genau ist folgendes Verfahren: Man ordne die sämmtlichen \(m\) Beobachtungsfehler (absolut genommen) nach ihrer Grösse, und nenne den mittelsten, wenn ihre Zahl ungerade ist, oder das arithmetische Mittel der zwei mittelsten bei gerader Anzahl, \(M\). Es lässt sich zeigen, was aber an diesem Orte nicht weiter ausgeführt werden kann, dass bei einer grossen Anzahl von Beobachtungen \(r\) der wahrscheinlichste Werth von \(M\) ist, und dass die wahrscheinlichen Grenzen von \(M\)
\[
r\bigl(1-e^{\rho\rho}\surd\tfrac\pi{8m}\bigr)\quad\text{und}\quad r\bigl(1+e^{\rho\rho}\surd\tfrac\pi{8m}\bigr)
\]
sind, oder die wahrscheinlichen Grenzen des Werthes von \(r\)
